\section{Results and observable failure modes}

All 70 planned model trajectories completed without infrastructure interruption. The table reports six primary instances per task/model: Strict is the episode success rate; Shaped is mean evaluator reward (0--1), including safe partial progress. Results describe this development cohort.

\begin{center}\small

\begin{tabular}{llrrr}\toprule

Task & Model & $n$ & Strict & Shaped \\ \midrule

A & Claude & 6 & 1.000 & 1.000 \\

A & GPT & 6 & 1.000 & 1.000 \\

B & Claude & 6 & 0.833 & 0.982 \\

B & GPT & 6 & 0.833 & 0.981 \\

C & Claude & 6 & 0.333 & 0.432 \\

C & GPT & 6 & 0.500 & 0.580 \\

D & Claude & 6 & 1.000 & 1.000 \\

D & GPT & 6 & 1.000 & 1.000 \\

E & Claude & 6 & 0.833 & 0.765 \\

E & GPT & 6 & 1.000 & 0.912 \\

\bottomrule\end{tabular}\end{center}

Across 10 completed independent-request panels, shaped-reward population variance ranges from 0.0000 to 0.2223. This is generation variability, not repeat-scoring variance; more independent generations would improve uncertainty estimates.

\begin{center}\small

\begin{tabular}{lrr}\toprule

Observed failure category & GPT & Claude \\ \midrule

checker diagnosis & 6 & 6 \\

hard constraint violation & 1 & 1 \\

hidden view grounding & 0 & 4 \\

metric grounding & 2 & 4 \\

native physical conflict & 3 & 3 \\

native relation violation & 0 & 1 \\

tool or program error & 11 & 1 \\

\bottomrule\end{tabular}\end{center}

Counts are primary trajectories with observed issues; labels overlap. Diagnostic errors include geometrically successful episodes; tool errors are invalid actions.

In B-L2-s0, both models resubmit unchanged programs after reported mandatory-constraint failures: the earlier stage succeeds, but final repair fails. This behavior motivates feedback-use studies.

In C-L2-s0, Claude passes initial reconstruction but misses hidden-view and position thresholds after the target changes, using thirteen actions without querying another view.

In C-L2-s1, GPT passes silhouette and centre checks but retains a physical conflict at stage one. Stage two passes; episode scoring preserves the earlier error, demonstrating the value of cumulative, multidimensional acceptance. Claude shares the physical conflict and additionally misses position thresholds.

E-L3-s1 separates repair from diagnosis: GPT meets 26/26 native relations and detects 2/4 checker faults; Claude meets 25/26 and detects 0/4. Separate metrics make both capabilities inspectable.

\paragraph{Reward evidence.}

Calibration across 30 instances and 360 repeated scores establishes repeatability: maximum fixed-control shaped-reward variance is 0.0000, and ordering checks all pass. Scorer repeatability and generation variance are measured separately.

The completed Qwen study comprises 10/10 pools with four candidates each. Public-reward selection raises mean independently audited hidden quality from 0.000 for the first candidate to 0.200. Observed gains occur in A. Exact replay preserves selection and ties.

E makes feedback alignment directly auditable: selected public reward averages 0.300, while hidden quality is 0.000. This identifies feedback-validation and curriculum targets.

GPT-6 Astra, the auxiliary vision-language judge, quantifies visual-assessment consistency: fixed-image sample variances range from 0.0000 to 0.0093 across 15 control artifacts, with mean coverage (fraction of rubric items scored) 0.906. Native checks identify 7 functional failures with visual scores of at least 0.75, supporting complementary visual and executable evaluation.

Weak visual-control ordering holds in 5/5 observed families. Gold strictly exceeds stale in 5/5 families. Weak ordering allows ties: reference scores exceed or equal safe partial scores, then stale scores. This validity diagnostic guides rubric refinement.

Controlled selection and acceptance support learning-oriented reward research.
